% =====================================================================
% FUNDAMENTAÇÃO TEÓRICA (Referencial Teórico)
% Reúna os conceitos necessários para o leitor compreender o trabalho.
% Este capítulo demonstra: figura, equação, citação e listas.
% =====================================================================

Lorem ipsum dolor sit amet, consectetur adipiscing elit, sed do eiusmod tempor
incididunt ut labore et dolore magna aliqua~\parencite{exemplo-livro}. Ut enim ad
minim veniam, quis nostrud exercitation ullamco laboris.

%-------------------------------------------------------------
\section{Primeiro conceito}
\label{sec:conceito-um}
%-------------------------------------------------------------

Duis aute irure dolor in reprehenderit in voluptate velit esse cillum dolore eu
fugiat nulla pariatur~\parencite{exemplo-artigo}. A \autoref{fig:exemplo} ilustra
o conceito.

% Citação DIRETA curta (até 3 linhas): fica no corpo do texto, entre aspas, com a
% fonte e a página logo após (NBR 10520:2023). Use \enquote{...} (csquotes) para as
% aspas — NUNCA a aspa reta ", que sob XeLaTeX sai como ”...” (fecha nas duas pontas).
Nesse sentido, \enquote{a citação direta de até três linhas permanece no corpo do
texto, entre aspas}~\parencite[p. 15]{exemplo-livro}.

% Citação DIRETA LONGA (mais de 3 linhas): destacada em parágrafo próprio, com
% recuo de 4 cm da margem esquerda, fonte menor, espaço simples e SEM aspas
% (NBR 10520:2023; guia UNIVALI). Use o ambiente \begin{citacao}...\end{citacao}:
% o abnTeX2 já aplica recuo (4 cm), \ABNTEXfontereduzida e espaçamento simples.
% A chamada (autor, ano, página) precede o bloco; não se usam aspas.
Já a citação com mais de três linhas é destacada do texto, como em
\textcite[p. 86]{exemplo-livro}:

\begin{citacao}
  Lorem ipsum dolor sit amet, consectetur adipiscing elit, sed do eiusmod tempor
  incididunt ut labore et dolore magna aliqua. Ut enim ad minim veniam, quis
  nostrud exercitation ullamco laboris nisi ut aliquip ex ea commodo consequat.
  Duis aute irure dolor in reprehenderit in voluptate velit esse cillum dolore eu
  fugiat nulla pariatur.
\end{citacao}

% --- EXEMPLO DE FIGURA -----------------------------------------------
% Para usar uma imagem real, troque o conteúdo do \centering pela linha:
%   \includegraphics[width=0.75\textwidth]{imagens/sua-figura.png}
% O comando \figuraplaceholder (definido em tex/_placeholders.tex) apenas
% desenha uma caixa cinza para o template compilar sem imagens externas.
\begin{figure}[!htb]
  \caption{Legenda da figura (acima da figura, conforme NBR 14724).}
  \label{fig:exemplo}
  \centering
  % A fonte fica alinhada à margem esquerda DA FIGURA (não do texto).
  \elementocomfonte{\figuraplaceholder{0.6\textwidth}{4cm}{imagens/sua-figura.png}}%
    {Fonte: Elaborada pelo autor (\the\year{}).}
\end{figure}

%-------------------------------------------------------------
\section{Formulação matemática}
\label{sec:formulacao}
%-------------------------------------------------------------

Equações são numeradas e podem ser referenciadas, como mostra a
\autoref{eq:exemplo}:

\begin{equation}
  E = m c^{2},
  \label{eq:exemplo}
\end{equation}

\noindent onde $E$ é a energia, $m$ a massa e $c$ a velocidade da luz no vácuo.
At vero eos et accusamus et iusto odio dignissimos ducimus.

%-------------------------------------------------------------
\section{Síntese}
%-------------------------------------------------------------

Os principais pontos abordados foram:

\begin{enumerate}[label=(\roman*)]
  \item primeiro ponto: lorem ipsum dolor sit amet;
  \item segundo ponto: consectetur adipiscing elit;
  \item terceiro ponto: sed do eiusmod tempor incididunt.
\end{enumerate}
